\documentclass[%
        paper=a4,
	% BCOR=10mm, % Binding CORrection
	oneside,			% two sided printing
        %cleardoublepage=empty,         % empty backsides for two sided printing
        twocolumn=false,
	titlepage=false,			% support for title environment
        abstract,
	numbers=noenddot,		% numbers (in headlines) without ending points
	headsepline,			% line below header
	%footsepline,			% line above footer
	headings=normal,                % size of header is normal
	fontsize=11pt,			% fontsize in points
	pagesize=pdftex,		% define typearea for pdftex usage
        headheight=46.96317pt,          % value based on MyThOS logo
        parskip=off                     % or half*
]{scrartcl}

% ------------------------------------------------------------------------------
% character encoding and fonts for the document
\usepackage[T1]{fontenc} % should be loaded before inputenc
\usepackage[utf8]{inputenc}	    % unicode, newer than utf8x
\usepackage{ae}
\usepackage{aecompl}
\usepackage[T1]{url}			% web adresses with t1 style
\urlstyle{tt}					% web adresses in tt-style
%\usepackage[UKenglish]{babel}
\usepackage[ngerman,UKenglish]{babel}

% ------------------------------------------------------------------------------
% font setup
\usepackage{microtype}
\usepackage{ellipsis}			% corrects wrong spaces after dots
\usepackage{amsmath}
%\usepackage{amsfonts}
\usepackage{amssymb}
\usepackage{latexsym}
\usepackage{dsfont}
\usepackage{textcomp}

%% bera mit euler
\usepackage{bera,amssymb,eufrak}
\usepackage[euler-digits,euler-hat-accent]{eulervm}

\setlength{\parskip}{0.5em plus 1em minus 0.5em}
\clubpenalty 2000
\widowpenalty 2000
%\setlength{\tolerance}{2000}
\setlength{\emergencystretch}{20pt}   % kann man evtl. auf 20 erhöhen
%\setlength{\hfuzz}{1pt}
\setlength{\hfuzz}{10pt}

% ------------------------------------------------------------------------------
% pdf creation
\usepackage[%
  pdftex,
  colorlinks=true,
  urlcolor=red,
  citecolor=red,
  linkcolor=red,
  bookmarks=true,
  breaklinks=true
]{hyperref}

% ------------------------------------------------------------------------------
% graphics
\usepackage{siunitx}
\sisetup{scientific-notation = true}%
\usepackage[usenames,dvipsnames]{xcolor}
\usepackage{graphicx}
\usepackage{subfigure}
\usepackage{tikz}
\usetikzlibrary{calc,trees,positioning,arrows,chains,shapes.geometric,%
    decorations.pathreplacing,decorations.pathmorphing,shapes,%
    matrix,shapes.symbols}
\usepackage{pgfplots}
\pgfplotsset{compat=newest}
\usetikzlibrary{pgfplots.statistics} % For Box plots
\usepgfplotslibrary{units} % Allows to enter the units nicely


% ------------------------------------------------------------------------------
% listings
\usepackage{listings}		% include listings
\newcommand{\code}{\lstinline}
\lstloadlanguages{C++}
\lstset{
  basicstyle=\ttfamily\footnotesize,
  keywordstyle=\bfseries,
  stringstyle=\textit,
  emphstyle=\textit,
  commentstyle=\textit,
  breaklines=true,
  breakatwhitespace=true,
  breakautoindent=true,
  morekeywords={},
  numbers=left,
  numberstyle=\tiny\color{black},
  stepnumber=1,
  numbersep=5pt,
  tabsize=2,
  showspaces=false,
  showstringspaces=false,
    xleftmargin=17pt, 			% left margin where linenumbers and frames do start
    framexleftmargin=17pt, 		% left margin of frame (which will be colored)
    framesep=0pt,
    belowcaptionskip=0pt,
    captionpos=b,
    framerule=1pt
}

\usepackage{scrhack}
\usepackage{float}
\floatstyle{plain}
\newfloat{listing}{tbhp}{myc}
\floatname{listing}{Listing}

% ------------------------------------------------------------------------------
% tables
\usepackage{multirow}
%\usepackage{ragged2e}
\usepackage{tabularx}
\usepackage{longtable}
\usepackage{booktabs}   % toprule, midrule in Tabellen
\newcolumntype{Y}{>{\raggedright\arraybackslash}X}
\usepackage{caption}

\usepackage{collcell}
\newcolumntype{L}{>{\collectcell\lstinline}l<{\endcollectcell}}
\newcolumntype{R}{>{\collectcell\lstinline}r<{\endcollectcell}}

\usepackage{bytefield}

%------------------------------------------------------------------------------
\usepackage{xspace}
\newcommand{\mythos}{\textsc{MyThOS}\xspace}
\def\cc{{C\nolinebreak[4]\hspace{-.05em}\raisebox{.4ex}{\tiny\bf ++}}}

\usepackage[textsize=tiny]{todonotes}

%\usepackage{titling}
\usepackage{enumerate}
\usepackage{epigraph}
\renewcommand{\epigraphsize}{\footnotesize}
\setlength{\epigraphwidth}{0.618\textwidth}
\usepackage{placeins}

\titlehead{
\begin{center}
  \noindent\rule{1.0\textwidth}{1pt}\\[1cm]
  \includegraphics[scale=0.5]{MyThOS_Logo}\\[0.5cm]
  \sffamily
  \small
  \begin{tabularx}{\textwidth}{lX}
    Förderkennzeichen: & 01IH130003 \\
    Vorhabensbezeichnung: & MyThOS \\
    & Modulares Betriebssystem für Massiv Parallele Anwendungen
  \end{tabularx}\\[0.5cm]
  \noindent\rule{1.0\textwidth}{1pt}
\end{center}
%\vspace{1cm}
}


\usepackage{scrlayer-scrpage}
\makeatletter
\title{\@myTitle}
\renewcommand*{\title}[1]{\def\@myTitle{#1}}
\setkomafont{pageheadfoot}{\sffamily\small}
\lohead{%
\@myTitle\\%
\headmark%
}
\cohead{}
\rohead{\includegraphics[height=1.5cm]{MyThOS_Logo}}
\makeatother
\pagestyle{scrheadings}

%\usepackage{showframe}

% ------------------------------------------------------------------------------
% layout
\usepackage{multicol}
\usepackage{marginnote}
\renewcommand*{\marginfont}{\footnotesize}
\typearea[current]{last}		% recalculate type area


\title{MyThOS D3.3 Softwareentwicklungsplan}
\author{Stefan Bonfert, Vladimir Nikolov, Robert Kuban, Randolf Rotta}

\hypersetup{
  pdftitle={MyThOS D3.3 Softwarentwicklungsplan},
  pdfsubject={MyThOS Deliverable},
  pdfauthor={Stefan Bonfert, Vladimir Nikolov, Robert Kuban, Randolf Rotta},
  pdfkeywords={MyThOS} % comma separated
}

\begin{document}
\selectlanguage{ngerman}
\maketitle

\begin{abstract}
Dieses Dokument umfasst den Strategieplan für die Weiterentwicklung des Prototypen und Anleitungen für Systementwickler. Abschnitt~\ref{sec:kernel-objects} gibt Beispiele und Hinweise für die Entwicklung eigener Kernel-Objekte, die Behandlung von Systemaufrufen, die Synchronisation asynchroner Aufgaben, sowie die Verwaltung von Beziehungen zwischen Kernel-Objekten. 
Abschnitt~\ref{sec:extensions} fasst nützliche Erweiterungen des bestehenden Kernels zusammen. Dazu gehören zum Beispiel Verbesserungen bei der Fehler- und Performanceanalyse, Schnittstellen zur Konfiguration der Interruptbehandlung und Capability Transfers zwischen Kommunikationsportalen. 
\end{abstract}

\newpage
\tableofcontents

% --- content ------------------------------------------------------------------
%\section{Übersicht}


% --- content ------------------------------------------------------------------
\selectlanguage{UKenglish}
\section{Kernel Development}
\label{sec:kernel-objects}

This section describes the steps necessary to create a new kernel
object. First, the \texttt{IKernelObject} interface
(Listing~\ref{lst:ikernelobject}) is discussed. It has to be
implemented by all kernel objects in order to participate in the
resource management and processing of system calls.  Then, the
allocation and initialisation of kernel objects through factories is
described.  Finally, an overview over the synchronisation monitors and
the weak references mechanism is given.

The kernel source includes an example object, which serves as
a starting point for the creation of custom Kernel Objects and as a
mock for testing object allocation and deletion. For further
information, consider the documentation embedded into the source
code. Most of the code embedded here is part of the Example object.

\begin{lstlisting}[float, label=lst:ikernelobject, caption=The \texttt{IKernelObject} interface.]
class IKernelObject : public ICastable {
public:
  virtual Range<uintptr_t> addressRange(Cap self);
  virtual optional<Cap> mint(Cap self, CapRequest request);
  virtual optional<void> deleteCap(Cap self, IDeleter& del) = 0;
  virtual void deleteObject(Tasklet* t, IResult<void>* r) = 0;
  virtual Range<uintptr_t> objectRange() const = 0;
  virtual optional<CapEntryRef> lookup(Cap self, CapPtr needle, CapPtrDepth maxDepth);
  virtual void invoke(Tasklet* t, Cap self, IInvocation* msg);
};

class ICastable {
public:
  virtual optional<void const*> vcast(TypeId id) const;
};
\end{lstlisting}

\subsection{Ressource management}

The \texttt{addressRange()} method is used by the resource inheritance
tree for parent-ship tests. The method has to return the physical
address range that is used by the object. The range has to be a
superset of the address ranges of all derived child objects in the
capability inheritance tree. Kernel objects that will have no children
can simply use the default implementation.

\texttt{deleteCap()} and \texttt{deleteObject()} are both part of the
multistep deletion process. \texttt{deleteCap()} notifies the object
that a capability is going to be deleted. For most objects, this will
have no effect for reference and derived capabilities. For original
capabilities, the objects recursively deletes its capability entry and
schedules itself for asynchronous deletion using the \texttt{IDeleter}
object given as a parameter. Listing~\ref{lst:deletecap} shows a simple example
for an object without any embedded capability entries.

\begin{lstlisting}[float, label=lst:deletecap, caption=Example handler for capability revocation.]
optional<void> ExampleObj::deleteCap(Cap self, IDeleter& del) {
  if (self.isOriginal()) del.deleteObject(del_handle);
  return Error::SUCCESS;
}
\end{lstlisting}

\texttt{deleteObject()} carries out the final step of the asynchronous
deletion. The object will use its deletion monitor to wait for all
outstanding asynchronous request to complete and then asynchronously
requests the deletion from the \texttt{UntypedMemory} object it was
allocated from. An example is shown in
Listing~\ref{lst:deleteobject}. The memory's \texttt{release()} method
will call the object's \texttt{objectRange()} to query the address
range that is used by the object itself. First, the objects virtual
destructor \texttt{~IKernelObject()} is called in order to let the
object free all contained additional memory blocks. Then, the objects
memory is freed. Finally, the deleter is notified about the completion
by replying to the \texttt{IResult} pointer.

\begin{lstlisting}[float, label=lst:deleteobject, caption=Example handler for asynchronous object deletion.]
void ExampleObj::deleteObject(Tasklet* t, IResult<void>* r) {
  monitor.doDelete(t, [=](Tasklet* t){ this->_mem->release(t, r, this); });
}
\end{lstlisting}

The \texttt{mint()} method is used when creating a reference or derived capability. This mechanism allows to restrict access rights, add communication badges and similar by altering the data portion of the capability according to the request argument. The meaning of the request arguments is defined by the actual object because the capability data portion is interpreted only by the object itself. If a kernel object does not accept the minting request, it should return an error code.

\subsection{System Calls as Capability Invocation}

The \texttt{lookup()} method is used by the system call entry code in
order to find the target object that matches the given capability
pointer. The search in the caller's capability space recursively walks
through its capability maps. The default implementation of this method
fails and returns an error code. A custom implementation is only
necessary if the object acts as a capability map.

The \texttt{invoke()} method implements the receiving side of the
capability invocation to a kernel object. Its default implementation
simply returns with an error. Listing~\ref{lst:invokehandler} shows an
example of the invocation handling. The \texttt{invoke()} method
dispatches the message based on the message label and asynchronously
calls the appropriate implementation. The implementation checks the
caller's access rights by inspecting the capability \texttt{self} that
was used to access the object. Then, the arguments are copied from the
message into local buffers and checked for validity. The copy is
crucial for security because otherwise concurrently running
application threads could manipulate the arguments after the sanity
check have been passed. After carrying out the request, the object
replies and releases the its monitor.

\begin{lstlisting}[float, label=lst:invokehandler, caption=Example invocation 
handler.]
void ExampleObj::invoke(Tasklet* t, Cap self, IInvocation* msg) {
  switch (msg->getLabel()) {
  case 0:
    return monitor.request(t, [=](Tasklet* t) { printMessage(t, self, msg); });
  default:
    msg->replyResponse(Error::INVALID_REQUEST);
  }
}

void ExampleObj::printMessage(Tasklet* t, Cap self, IInvocation* msg) {
  // may use data portion of self for access control and badges...
  auto msgdata = msg->getMessage();
  // do something useful with the message data...
  message.replyResponse(Error::SUCCESS);
  monitor.requestDone();
}
\end{lstlisting}

The \texttt{IInvocation} interface provides several more methods that
help, for example, to look up additional kernel objects in the callers
capability space, access the callers \texttt{ExecutionContext}, and
the caller's logical address space configuration.

\subsection{Interactions between Kernel Objects}

When a capability lookup is used to retrieve a pointer to another
kernel object, this pointer is still typed as
\texttt{IKernelObject}. This is sufficient to pass invocation messages
to the object. However, the invocation mechanism is complex and
possibly costly. Therefore it would be more efficient and convenient
to call the kernel object's implementation-specific methods
directly.

Normal C++ programs use the compiler-generated runtime type
information (RTTI) and the \texttt{dynamic\_cast} conversion. This
information is not available in the kernel in order to reduce its
size. Instead, the conversion is accomplished through the
\texttt{ICastable} mechanism. The method \texttt{vcast()} is
implemented by each kernel object in order to hand out pointers to
more specific interfaces. Listing~\ref{lst:vcast} shows an example.
The converted pointer is then retrieved via the \texttt{cast()} helper
method as shown in Listing~\ref{lst:cast}.

\begin{lstlisting}[float, label=lst:vcast, caption=Example run-time type conversion in the \texttt{UntypedMemory} object.]
virtual optional<void const*> vcast(TypeId id) const override {
  if (id == TypeId::id<UntypedMemory>()) return this;
  if (id == TypeId::id<IAllocator>()) return static_cast<IAllocator const*>(this);
  return Error::TYPE_MISMATCH;
}
\end{lstlisting}

\begin{lstlisting}[float, label=lst:cast, caption=Casting a kernel object into a more specific type.]
auto alloc = ptr->cast<IAllocator>();
if (alloc) { // checks whether the conversion was successful
  alloc->allocate(...);
}
\end{lstlisting}



\subsection{Object Creation}

Kernel objects are created with memory from a single
\texttt{UntypedMemory} via the help of a factory object.  The only
exception are a few initial objects that exist statically.  The
factory has to implement the \texttt{IFactory} interface, which is
shown in Listing~\ref{lst:ifactory}. 

\begin{lstlisting}[float, label=lst:ifactory, caption=The \texttt{IFactory} interface.]
class IFactory {
public:
  virtual optional<void> factory(Cap umcap, IAllocator* mem, IInvocation* msg,
                                 CapEntry* tgtentry) = 0;
};
\end{lstlisting}

The \texttt{umcap} capability belongs to the \texttt{UntypedMemory}
that called the factory.  This value is later needed to complete the
insertion into the resource tree.  The \texttt{mem} argument points to
the \texttt{IAllocator} interface of the \texttt{UntypedMemory}.
Because the memory calls the factory the memory's synchronisation
monitor has been entered, synchronous calls to \texttt{mem} can be
used safely.  The created object has to store \texttt{mem} in order to
be able to release its memory later.  The \texttt{tgtentry} points to
the capability entry that shall receive the first capability that
points to the newly created kernel object.  Finally, the \texttt{msg}
arguments contains the invocation buffer that was passed from userland
with additional initialisation parameters for the new object.

An example factory implementation is shown in
Listing~\ref{lst:factory}.  First, the target capability entry is
acquired and locked in order to prevent data races (line 3).  Then,
memory for the object is allocated and the object is constructed in
this memory.  The factory and the object constructor can allocate
additional memory from \texttt{mem} (line 2).  If the allocation
fails, the memory has to be released and the target capability entry
has to be reset (lines 5--7).

Finally, the original capability of the created object is inserted as
child of the \texttt{UntypedMemory} that it was allocated from (line
10).  The parent entry is retrieved from the invocation message and
should match \texttt{umcap}.  The operation fails when the parent
object was deleted concurrently.  In this case, the object and the
target entry have to be released (lines 11--14).  The original
capability can be either stored as member variable in the created
object, as shown in the example, or can be stored directly in the
target entry.  In the example, a reference capability is written to the
target entry.

\begin{lstlisting}[float, label=lst:factory, caption=An example factory implementation.]
optional<void> ExampleFactory::factory(Cap umcap, IAllocator* mem,
                                       IInvocation* msg, CapEntry* tgtentry) {
  if (!tgtentry->acquire()) return Error::LOST_RACE;
  auto obj = mem->create<ExampleObj,64>(mem, ...); // create with 64b alignment
  if (!obj) {
    tgtentry->reset();
    return Error::INSUFFICIENT_RESOURCES;
  }
  Cap tgtcap(obj);
  auto res = cap::inherit(msg->getCapEntry(), obj->ownRoot, umcap, tgtcap);
  if (!res) {
    mem->release(obj);
    tgtentry->reset();
    return res.state();
  }
  return cap::inherit(obj->ownRoot, tgtentry, tgtcap, tgtcap.asReference());
}
\end{lstlisting}



\subsection{Serialisation Monitors}

The monitors are used by kernel objects to serialise asynchronous
method calls and select the place (hardware thread) that processes the
calls.  Currently, four monitor variants are available.  The choice
depends on the synchronisation needs.  In most cases, the
\texttt{NestedMonitorDelegating} is a good choice.

\paragraph{DeletionMonitor.}
This monitor implements a reference counter that is used to delay the
deletion request until all counted references were released.  It is
used as base class for the other monitors.  The methods
\texttt{acquireRef()} and \texttt{releaseRef()} increase or decrease
the reference counter, respectively.  The \texttt{doDelete()} method
schedules a Tasklet to be executed the next time the reference counter
reaches zero.  The Tasklet is processed at the place of the
\texttt{releaseRef()} caller.  If the reference counter is zero
already, the Tasklet processed immediately.

\paragraph{SimpleMonitorHome.}
This variant comes close to the classic monitor concept.  The
\texttt{request()} method schedules the given Tasklet at a predefined
place called \texttt{home}.  At the end of each request handler
implementation, \texttt{requestDone()} has to be called in order to
release the reference counter.  Because all Tasklet scheduled by the
monitor are processed on the same hardware thread, they all are
mutually exclusive.

\paragraph{NestedMonitorHome.}
This extension of the previous monitor differentiates between
asynchronous \emph{requests} and \emph{responses} in order to
implement a form of nested locking.  The next requests is processed
only after the previous has finished by calling
\texttt{requestDone()}.  Requests can issue sub-requests to other
kernel objects and the last response has to call
\texttt{responseAndRequestDone()} instead of just
\texttt{responseDone()}.  Developers should be careful with cyclic
dependencies between kernel objects because the mutually exclusive
request are prone to deadlocks -- like in any scenario with nested
locks that can be entered just once.

\paragraph{NestedMonitorDelegating.}
This monitor has the same \emph{request}/\emph{response} interface as \texttt{NestedMonitorHome} but
the monitor is not bound to a specific home place.  Instead, the first
request that acquires exclusive access sets the home to its caller's
place.  All responses to sub-requests will processed at this
place.  Further requests that do not acquire exclusive access will be
processed at this place, too.  The last \texttt{requestDone()} or
\texttt{responseAndRequestDone()} releases the exclusive access, such
that later requests select a new home.  This behaviour is based on the
concept of \emph{delegation locks}.


\subsection{Connecting Objects through Weak Reference Capabilities}

Pointers between kernel objects require special care because the
target object can be deleted concurrently.  This would leave dangling
pointers, which result in fancy bugs and headaches.  In many cases,
the associations between kernel objects are established by the
transfer of access rights.  In \mythos this is implemented by passing
a capability to the kernel object and the capability revocation can be
used to clean up such pointers.  This idea is implemented by the
\texttt{CapRef} class.

\begin{lstlisting}[float, label=lst:capref, caption=The weak reference capability interface.]
template<class Subject, class Object, class Revoker=RevokeByUnbind>
class CapRef : public CapRefBase {
public:
  optional<void> set(Subject* subject, CapEntry& src, Cap srcCap);
  optional<Object*> get() const;
  void reset();
protected:
  virtual void revoked(Cap self, Cap orig);
};   
\end{lstlisting}

Listing~\ref{lst:capref} summarises the principal interface of
\texttt{CapRef}. Basically, the object contains a capability entry
that holds a reference capability and a pointer to the
\texttt{subject} capability holder.  The \texttt{set()} method
inherits the reference from the \texttt{src} and \texttt{srcCap} pair.
It fails if a concurrent call to \texttt{set()} was faster or when the
source entry was deleted concurrently.  The \texttt{get()} method
returns a pointer to the capability's object casted to the desired
object interface.  It fails if the capability was never set or was
revoked meanwhile.  The \texttt{reset()} method clears the reference
and called automatically from \texttt{set()}.

During capability revocation, the \texttt{revoked()} method will be
called synchronously.  It can be overridden by derived types.  The
default implementation calls 
\lstinline{Revoker::apply(subject, object, self, orig)}.
The default \texttt{RevokeByUnbind} calls
\lstinline{subject->unbind(object)} synchronously in order to notify the reference
holder. An example usage is given in Listing~\ref{lst:caprefex}.

\begin{lstlisting}[float, label=lst:caprefex, caption=An example use of weak references.]
class Foo : public IKernelObject {
public:
  optional<void> bindPortal(CapEntry& pcap); // calls portal.set(...)
protected:
  CapRef<Foo, IPortal> portal;
  friend class RevokeByUnbind;
  void unbind(optional<IPortal*> o); // reacts to the revocation
};
\end{lstlisting}


\section{Strategy for further Development}
\label{sec:extensions}

\subsection{Capability Transfer}

Capability transfer and capability unwrapping in inter-process calls
are not necessary for the applications targeted by \mythos.  However,
they promise to simplify the interaction between unrelated application
without a common supervisor.

In order to support the future implementation of such transfers, the
invocation buffer format mirrors the seL4 message format and reserves
the necessary space.  In order to support a capability transfer, a
portal implementation would read the desired target capability address
from the receiver's invocation buffer, look up the respective entry in
the receiver's capability space, and insert a reference capability
there.

\subsection{Image-based Debugging Facilities}

Because the system memory in \mythos is just 4GiB large, it is
feasible to create an image of the whole system memory by reading it
from the host over PCIe.  Host access to the coprocessor memory is
cache coherent if the host-side caching is disabled. Although it is
possible create a kernel image without further kernel support,
concurrent changes to the kernel memory can easily lead to an
inconsistent image and interesting details about the state of each
hardware thread are hidden in their registers.

Two approaches can be implemented in order to obtain a consistent
kernel image: \emph{Snapshot Images} capture a consistence state by
issuing a non-maskable interrupt (NMI), which dumps the processor
state into each thread's NMI stack and, and then, busy waits in the
interrupt handler until the image is captured.  In contrast, a
\emph{Checkpoint Image} is created by using a interrupt to force all
threads into the kernel, but instead of waiting in the NMI handler,
all threads wait before the next Tasklet is executed. This allows to
create a kernel image that is easier to analyse because all
transactions and critical sections were completed.

\subsubsection{Doorbell Interrupts and User-Level Interrupt Handling}
The Intel Xeon Phi Processor has a doorbell interrupt, which can be
triggered from the host software in order to wakeup the sleeping
processor or interrupt its current activity in order to handle urgent
requests. This mechanism is based on configuration registers in the
processor's PCIe MMIO space and the XeonPhy's IOMMU.  The respective
configuration can be added through a KNC-specific source code
module. One imminent application would be the interrupt broadcast for
debugging purposes.

The present interrupt handling in \mythos is sufficient to handle
traps and interrupts inside the kernel but is not extensible. A useful
improvement would be the implementation of user-level interrupt
handling. The basic idea is to represent interrupt gates as kernel
objects and triggered interrupts send a message or notification to a
portal or execution context.

\subsection{Exploiting Hardware Transactional Memory}
% was ist das / wozu?
Hardware Transactional Memory such as Intel's TSX extension can speed
up the execution by enabling lock elision and much simpler
non-blocking algorithms.  Previous related work about transactional
memory in
microkernels~\cite{smejkal2015transactional,fuchs2014hardware} focused
on the latency of the IPC path.  The results were disappointing
because one or two atomic exchange or CAS operations are difficult to
beat.

% status in mythos?
% wie könnte es besser unterstützt werden?
\mythos contains quite a lot non-blocking algorithms, especially for
double-linked queues of pending messages and the resource inheritance
tree. These can be replaced by hardware transactional memory in order
to reduce complexity and busy waiting. Very likely, TSX would not
speed up the serialisation monitors and the tasklet queue.

\subsection{Conservative Sleeping}

When there is no work to do, a hardware thread should switch into a
sleep state in order to minimize energy consumption and to lend its
thermal budget to other cores.  Currently, a hardware thread in
\mythos starts sleeping as soon as there are no more executable
Tasklets and Execution Contexts (application threads).  However, this
is not optimal because waking up from a sleep state requires an
interrupt to be send by another thread, which imposes a latency
penalty. Entering and exiting deeper sleep modes also adds a
considerable overhead.

A possible solution to this problem is introducing an active waiting
phase, which resembles the \emph{pause} phase for mutexes first
introduced by Ousterhood.  Heuristics might be used to find a tradeoff
between active waiting phase length, hence energy consumption, and the
latency penalty of waking up from a sleep state. \mythos can support
the development of such heuristics by providing diagnostic data like
active waiting time or sleeping phase length.

\subsection{Sleeping through MWAIT}

Many x86 Prozessors, except the Intel XeonPhi Knight Corner, support a more efficient sleep\&wakeup mechanism for the hardware threads. It is based on monitoring a cache line for access by other hardware threads with the \texttt{MONITOR} instruction. Then the thread can sleep with \texttt{MWAIT} and is woken up whenever the monitored line is modified or evicted from the thread's cache. 

The benefit of this approach is, that client threads that send a Tasklet do not need to check whether they have to send an IPI interrupt. The sleeping thread does not need to go through the interrupt entry routine when waking up. In \mythos this would allow to leave out the release of the thread's Tasklet queue prior to sleeping.

\subsection{Adapt Tracing for Tasklet Queues and Synchronisation Monitors}
The previous implementation of \mythos used fixed-size Ringbuffers for the exchange of Tasklets between hardware threads. In addition, Tasklets were allowed to enter multiple queues at the same time and were allocated dynamically. This introduced interesting challenges for the global reconstruction of thread-local traces. 

In contrast, the present design is much simpler because all Tasklets have to follow a request-response cycle between kernel objects. This allows to simplify the trace recording and reconstruction. The port of the trace subsystem to the new design is still pending.


\bibliographystyle{alpha}
\bibliography{literature}

\end{document}
